\section{Proof attempt for the strongest candidate}

\subsection{History-based survival lemma}
\mstatus{PROVED} \istatus{PARTIAL} \nstatus{NOT CLAIMED} The lemma below is
a standard finite-horizon hitting bound,
included as the proof engine rather than as a novelty claim.

Let $Y_j$ be the indicator of the declared discovery event at comparison
opportunity $j$ and let $\tau=\inf\{j\geq1:Y_j=1\}$. The event
$\{\tau\geq j\}$ is determined by the history before opportunity $j$.
Assume there are deterministic numbers $p_j\in[0,1]$ such that, on every
pre-opportunity history with no previous hit,
\[
 \Prb(Y_j=1\mid\mathcal{F}_{j-1})\geq p_j.
\]
This is a history-uniform conditional hazard condition. It allows the target
attached to the event to vary with $j$. The index $j$ is deliberately not the
BQ cycle index $k$: it may represent a serialized mutation attempt, a full
generation event $D_k^{\rm gen}$, or an acquisition event
$D_k^{\rm acq}$, provided the choice and filtration are fixed.

\begin{lemma}[Non-stationary survival and cost bound]
For every $n\geq1$,
\[
 \Prb(\tau>n)\leq\prod_{i=1}^{n}(1-p_i).
\]
If opportunity costs $c_i$ are deterministic, then
\[
 \Ex[C_\tau]
 \leq \sum_{i\geq1}c_i\prod_{j=1}^{i-1}(1-p_j),
\]
with the inequality valid also when the right side is infinite. For a budget $B$, assume
$\sum_{i=1}^\infty c_i=\infty$ and let
$n(B)=\max\{n\geq0: \sum_{i=1}^{n}c_i\leq B\}$, with empty sum zero.
\[
 \Prb(C_\tau\leq B)\geq
 1-\prod_{i=1}^{n(B)}(1-p_i).
\]
\end{lemma}

\begin{proof}
Write $S_i=\{\tau\geq i\}$ for survival to the start of opportunity $i$.
Since $S_i\in\mathcal{F}_{i-1}$,
\begin{align*}
 \Prb(S_{i+1})
 &=\Ex\!\left[\mathbf{1}_{S_i}
       \Prb(Y_i=0\mid\mathcal{F}_{i-1})\right] \\
 &\leq (1-p_i)\Prb(S_i).
\end{align*}
Starting from $\Prb(S_1)=1$ and iterating gives
$\Prb(S_{n+1})\leq\prod_{i=1}^{n}(1-p_i)$, which is the first claim.
For the cost, use the nonnegative-sum identity
\[
 C_\tau=\sum_{i\geq1}c_i\mathbf{1}_{\{\tau\geq i\}},
 \qquad
 \Ex[C_\tau]=\sum_{i\geq1}c_i\Prb(S_i),
\]
and apply the survival bound termwise. Finally, completion of any of the
first $n(B)$ opportunities costs at most $B$, so budget success is at least the
probability of a hit among those opportunities. The first claim supplies the
stated lower bound.
\end{proof}

\subsection{Deriving the window comparison}
For Candidate~1, apply the survival lemma on the $W$-opportunity stable-target
window. The memory policy satisfies $p_j=p_M$ and the comparator satisfies
an upper hazard bound $q_j^G\leq p_G$ on every no-hit history. The same
induction, with the inequality reversed for $G$, yields
\[
 \Prb(\tau_M>j)\leq(1-p_M)^j,
 \qquad
 \Prb(\tau_G>j)\geq(1-p_G)^j.
\]
Therefore, for $n\leq W$,
\[
 \Prb(\tau_M\leq n)\geq1-(1-p_M)^n
 >1-(1-p_G)^n\geq\Prb(\tau_G\leq n).
\]
For $2\leq n\leq W$, using
$\Ex[\min(\tau,n)]=\sum_{j=0}^{n-1}\Prb(\tau>j)$ and multiplying by the
common opportunity cost $c$ gives the strict truncated compute-cost comparison in
Proposition~\ref{prop:window}. At $n=1$ both policies pay the first opportunity
cost in the truncated-cost convention, so only the success-probability
comparison is strict there.

\subsection{Dependency audit}
\begin{itemize}
\item \textbf{Conditional independence:} not used. The proof requires
history-conditional lower/upper hazard bounds; dependence across offspring or
generations is allowed if those bounds remain true after every no-hit history.
\item \textbf{Markovianity:} not used in the hitting proof. It would only be
needed if one chooses a finite-state transition matrix to calculate the
hazards. A niche aggregation is valid for that purpose only under lumpability.
\item \textbf{Mutation bounds:} the quantity $\alpha$ must be a lower bound
on reaching the true current hard region conditional on the useful parent
being selected. A graph path's existence alone is insufficient; its
probability and compute cost matter.
\item \textbf{Parent retention and scheduling:} a lower bound $r$ must concern
availability of a genuinely useful parent at the declared opportunity.
Historical presence before a target shift does not establish continued
usefulness afterward. The S6-AE parent/operator schedule is deterministic:
conditional on $P_k$, each parent receives all 11 operators. Hence no random
parent-selection factor $\gamma$ may be inserted; a batch reachability bound
$\alpha_{\rm batch}$ must be derived under the declared RNG ensemble.
\item \textbf{Surrogate accuracy:} no accuracy is needed for the abstract
survival lemma. To pass from $D_k^{\rm gen}$ to
$D_k^{\rm acq}$, a retention condition connecting $\widehat h_k$ and the QD
rule to true $h_k^{\rm BQ}$ and future reachability is essential. Perfect
current-hardness prediction alone is insufficient when tied candidates have
different mutation futures. Without such a bridge, the implementation-level
step is OPEN.
\item \textbf{Cost normalization:} the comparison assumes equal deterministic
cost $c$ per opportunity. If BLADE archive updates, HGB40 inference, or BQ
validation change cost, opportunity-count dominance does not imply compute-cost
dominance.
\end{itemize}

\subsection{Failure boundary}
\mstatus{OPEN for a BLADE-specific dominance claim}
\istatus{PARTIAL} \nstatus{OPEN} The proof does not establish that BLADE's conditional hazard is
above the fixed-generator or HAC hazard. It fails to produce a BLADE result
at the exact step where one would need verified bounds on $r$ and
$\alpha_{\rm batch}$ (or a direct generation/acquisition hazard bound) and a
comparator upper bound $p_G$, all in the same compute units. The
counterexamples show these inequalities can reverse. No unproved transition
from a structural-memory narrative to hazard dominance is made.
